% !TeX root = ../教材题库_全部条目_转写工作稿.tex
% Source: XJ-B1-C05-EX-P182-Y1522
% 第5章 / 5.3 三角函数的图象与性质 / 教材 p.182、183 / PDF 第189、190页
% OCR未校对：公式与文字均可能有误；未完成数学TeX转写；禁止直接发给学生
% 原题及相邻原解截图（核对用）：../题目原图/XJ-B1-C05-EX-P182-Y1522-01.webp, ../题目原图/XJ-B1-C05-EX-P182-Y1522-02.webp
% 本文件只在转写完成后作为题干使用；OCR 条目尚不是可靠题干。
\begin{bankitem}{5.3 例题 例题3}{第5章 / 5.3 三角函数的图象与性质 / 教材 p.182、183 / PDF 第189、190页}{XJ-B1-C05-EX-P182-Y1522}{OCR未校对：公式与文字均可能有误；未完成数学TeX转写；禁止直接发给学生}
\begin{ocrdraft}
[教材p.182 / PDF第189页]
\par
0河利用三角函数的单调性，比较下列各组数的大小：
\par
11丌
\par
12丌
\par
（2）cos
\par
，COS
\par
7
\par
7
\par
[教材p.183 / PDF第190页]
\par
解（D由于一<一1.1<一1<，且3一，，在区间[
\par
丕〕上单调递增，
\par
因此
\par
sin(—l)>sin(—l.1)．
\par
11丌12丌
\par
（2）由于丌<一一<一一<2丌，且3，一cos在区间[丌，2丌]上单调递增，
\par
11丌
\par
12丌
\par
因此
\par
COS一一一一一<COS
\end{ocrdraft}
\end{bankitem}
